\documentclass[12pt,a4paper]{report}

\usepackage[utf8]{inputenc}
\usepackage[T1]{fontenc}
\usepackage[english]{babel}
\usepackage{array}
\usepackage{float}
\usepackage{geometry}
\usepackage{longtable}
\usepackage{listings}
\usepackage[hyphens]{url}
\usepackage[hidelinks,breaklinks=true]{hyperref}

\geometry{margin=2.5cm}
\setlength{\parindent}{0pt}
\setlength{\parskip}{0.6em}
\renewcommand{\arraystretch}{1.2}

\lstset{
  basicstyle=\small\ttfamily,
  breaklines=true,
  breakatwhitespace=false,
  columns=fullflexible,
  frame=single
}

\title{LifeRouter\\User Manual}
\author{Kentalives}
\date{June 2026}

\begin{document}

\pagenumbering{roman}
\maketitle
\setcounter{page}{2}
\tableofcontents
\clearpage
\pagenumbering{arabic}

\chapter{User Manual}

This manual contains the information needed to use the public API of LifeRouter.
 It is intended for developers who want to integrate the
service from Go or from another process that can communicate through NATS. It
does not describe the internal algorithmic details; those belong to the design
documentation and the technical manual.

\section{Prerequisites}

Using the service requires an accessible NATS process, a rasterizable map
configured as a grid, and external geotruth services that provide object
information. Public API coordinates are expressed in meters as
\texttt{[x, y, z]} triples. Conversion between meters and cells depends on
\texttt{grid.cellsizem}; by default, one cell represents \texttt{0.2} meters.

Agent identifiers received by the API must exist in geotruth. The service reads
the agent's real position, computes a route on the grid, and publishes movement
updates through the configured publishing dependency.

\section{Installation and Configuration}

This section describes how to start the service as a standalone process.
Embedded execution from another Go program is described later in the API usage
section.

\subsection{Environment Preparation}

The project is a Go module and requires a Go version compatible with the version
declared in \texttt{go.mod}. Besides the service code, the installation needs a
reachable NATS server, geotruth services that store and publish object
positions, and the map images referenced by the configuration.

Before starting the service, check these points:

\begin{itemize}
\item the NATS process is available at the configured URL;
\item the map image paths exist from the execution directory;
\item grid dimensions match the map that will be rasterized;
\item agent identifiers and signaling nodes exist in geotruth;
\item floor and portal configuration represents the real space.
\end{itemize}

\subsection{Build}

The normal build is run from the repository root:

\begin{lstlisting}
make build
\end{lstlisting}

This target runs \texttt{go build} with the default \texttt{all} build tag. The
Makefile keeps \texttt{MODULES} as an optional build-tag override for specific
runs, for example:

\begin{lstlisting}
make build MODULES="tag1 tag2"
\end{lstlisting}

Ordinary builds should use \texttt{make build} without overrides.

When Go dependencies change, update module metadata and the \texttt{vendor}
directory:

\begin{lstlisting}
make deps
\end{lstlisting}

\subsection{Running as a Service}

The executable loads configuration with \texttt{config.LoadConfig(nil)}, prepares
default dependencies, initializes the rasterized world, opens the service NATS
connection, and registers the dispatcher subjects. The default dependencies also
create internal query and publish pools toward geotruth. The process keeps
running until it receives a shutdown signal such as \texttt{SIGINT} or
\texttt{SIGTERM}. At that point it closes active agent sessions, stops the
emergency if one is running, drains the service NATS connection, and releases
closeable dependency resources.

After building, run the generated binary from a directory where
\texttt{config/config.yaml} and the configured map paths are accessible:

\begin{lstlisting}
./LifeRouter
\end{lstlisting}

The standalone entry point reuses \texttt{embedded.DefaultDependencies}. That
function creates separate internal NATS pools for geotruth queries and
publishes, and creates a simulated external system for part of the domain data.
For a real deployment, the integration should provide its own dependencies
through embedded execution or through a service that wraps this module.

\subsection{Configuration File}

The standalone service reads \texttt{config/config.yaml}. Omitted fields use
\texttt{DefaultConfig} values. Embedded calls, tools, or benchmarks can call
\texttt{LoadConfig} with an explicit path to use another YAML file, for example a
scenario configuration under \texttt{data/}. Keys can also be overridden with
environment variables that start with \texttt{PTHFSERVICE}; dots in key paths
become underscores. For example:

\begin{lstlisting}
PTHFSERVICE_APP_SERVERURL=nats://localhost:4222
PTHFSERVICE_GRID_CELLSIZEM=0.2
\end{lstlisting}

The following example shows the main part of a configuration with multiple
floors, rasterizable images, and portals:

\begin{lstlisting}
app:
  serverurl: "nats://localhost:4222"
  stderr: true

grid:
  cellsizem: 0.2
  rows: 63
  cols: 118
  pxpercell: 10
  wallauracells: 1
  layers:
    - name: "0"
      img: "./data/map.png"
    - name: "1"
      img: "./data/floor1.png"
  portals:
    - from: [1.2, 3.2]
      fromlayer: 0
      to: [1.2, 3.2]
      tolayer: 1
      traversalcost: 15
      bidirectional: true

pathfinding:
  freemoveheight: 2
  agent:
    debug: false
    findpathhandlerworkers: 8
    blockingwaithandlerworkers: 1024
    movefmetershandlerworkers: 16
  geotruth:
    queryconnections: 8
    publishconnections: 8
  emergency:
    debug: false
\end{lstlisting}

\subsection{Main Fields}

The \texttt{app} group defines NATS connectivity and log output. The
\texttt{serverurl} key contains the NATS server URL. The \texttt{stderr},
\texttt{fileout}, and \texttt{logfile} keys control whether logs are written to
standard error, to a file, or to both.

The \texttt{grid} group defines the rasterized representation of the space.
\texttt{rows} and \texttt{cols} are expressed in cells. \texttt{pxpercell}
indicates how many image pixels correspond to one grid cell.
\texttt{cellsizem} indicates how many meters one cell represents. Entries in
\texttt{layers} bind each floor to an image, and entries in \texttt{portals}
connect coordinates on different floors. Portal coordinates are expressed in
meters, not cells; portal layer fields are integer indexes into the
\texttt{layers} list.

The \texttt{pathfinding} group contains routing parameters:

\begin{itemize}
\item \texttt{freemoveheight}: free height used for direct movement and floor
queries;
\item diagnostic and concurrency parameters: debug traces and worker limits for
some agent NATS handlers;
\item emergency parameters: lighting hysteresis and preferred-route width;
\item agent parameters: vision radius and number of cells between real-position
updates;
\item geotruth parameters: number of internal NATS connections used for queries
and publishes toward the external service.
\end{itemize}

Tune these values together with cell size and map scale.

\section{Complete Configuration Reference}

This reference summarizes the fields accepted by the YAML configuration file,
either \texttt{config/config.yaml} or an explicit path. If a field is omitted,
the service uses the value from \texttt{DefaultConfig}. Environment variables
use the \texttt{PTHFSERVICE} prefix and replace dots with underscores. For
example, \texttt{grid.cellsizem} can be overridden with
\texttt{PTHFSERVICE\_GRID\_CELLSIZEM}.

\subsection{\texttt{app}}

\begin{description}
\item[\texttt{serverurl}] Text containing the NATS server URL. The default is
the local NATS default URL.
\item[\texttt{stderr}] Boolean. If \texttt{true}, logs are written to standard
error. The default is \texttt{true}.
\item[\texttt{fileout}] Boolean. If \texttt{true}, logs are written to a file.
The default is \texttt{false}.
\item[\texttt{logfile}] Text containing the log file path when
\texttt{fileout} is enabled. The default is an empty string.
\end{description}

\subsection{\texttt{grid}}

\begin{description}
\item[\texttt{layers}] List of floors. Each element has \texttt{name}, the
external floor or region identifier, and \texttt{img}, the path of the image to
rasterize. By default there is one floor named \texttt{0} with image
\texttt{../../data/map.png}.
\item[\texttt{rows}] Unsigned integer with the number of grid rows. The default
is \texttt{61}.
\item[\texttt{cols}] Unsigned integer with the number of grid columns. The
default is \texttt{114}.
\item[\texttt{pxpercell}] Decimal number of image pixels that correspond to one
cell. The default is \texttt{10}.
\item[\texttt{cellsizem}] Decimal number of meters represented by one cell. It
must be greater than \texttt{0}. The default is \texttt{0.2}.
\item[\texttt{wallauracells}] Integer number of margin cells around walls. It
must be greater than or equal to \texttt{0}. The default is \texttt{1}.
\end{description}

\subsection{\texttt{grid.portals}}

Each portal connects two coordinates on different floors. Coordinates are
expressed in meters as \texttt{[x, y]} lists. The service converts them to cells
using \texttt{grid.cellsizem}.

\begin{description}
\item[\texttt{from}] List \texttt{[x, y]} with the source coordinate in meters.
\item[\texttt{fromlayer}] Integer index of the source layer in
\texttt{grid.layers}.
\item[\texttt{to}] List \texttt{[x, y]} with the destination coordinate in
meters.
\item[\texttt{tolayer}] Integer index of the destination layer in
\texttt{grid.layers}.
\item[\texttt{traversalcost}] Integer traversal cost for crossing the portal.
\item[\texttt{bidirectional}] Boolean. If \texttt{true}, the portal can be used
in both directions. If \texttt{false}, it is used only from source to
destination.
\end{description}

No portal is defined by default.

\subsection{\texttt{pathfinding}}

\begin{description}
\item[\texttt{freemoveheight}] Decimal free height, in meters, used for direct
movement and floor queries. It must be greater than \texttt{0}. The default is
\texttt{2}.
\end{description}

\subsection{\texttt{pathfinding.emergency}}

\begin{description}
\item[\texttt{lighthystheresis}] Integer cost hysteresis applied to emergency
signaling. It must be greater than \texttt{0}. The default is \texttt{100}.
\item[\texttt{priopathwidth}] Decimal width, in grid cells, of emergency
preferred routes. It must be greater than \texttt{0}. The default is \texttt{1}.
\item[\texttt{debug}] Boolean. If \texttt{true}, enables diagnostic output for
emergency pathfinding. The default is \texttt{false}.
\end{description}

\subsection{\texttt{pathfinding.agent}}

\begin{description}
\item[\texttt{visionradiuscells}] Unsigned integer with the agent vision radius
in cells. It must be greater than \texttt{0}. The default is \texttt{10}.
\item[\texttt{cellsforrealupdate}] Integer number of cells the agent can move
before updating its real position. It must be greater than \texttt{0}. The
default is \texttt{5}.
\item[\texttt{maxreusableworldlinkedfloors}] Unsigned integer with the maximum
number of linked floors that can be reused in an agent world. The default is
\texttt{2}.
\item[\texttt{sparsevaluepagedirectory}] Boolean. Enables sparse storage for the
agent value directory. The default is \texttt{false}.
\item[\texttt{findpathhandlerworkers}] Integer maximum number of worker
goroutines that can handle \texttt{AgentFindPath} NATS requests in parallel. The
default is \texttt{GOMAXPROCS}.
\item[\texttt{blockingwaithandlerworkers}] Integer maximum number of worker
goroutines that can handle \texttt{BlockingWait} requests in parallel. The
default is \texttt{1024}.
\item[\texttt{movefmetershandlerworkers}] Integer maximum number of worker
goroutines that can handle \texttt{MoveFMeters} requests in parallel. The
default is \texttt{2 * GOMAXPROCS}.
\item[\texttt{debug}] Boolean. If \texttt{true}, enables diagnostic output for
agent pathfinding. The default is \texttt{false}.
\end{description}

These worker limits affect only the NATS handler layer. They do not change the
pathfinding algorithm and do not parallelize the steps of a single agent. Orders
for one communicator remain ordered by its internal protocol, while requests
for different agents can be handled concurrently.

\subsection{\texttt{pathfinding.geotruth}}

\begin{description}
\item[\texttt{queryconnections}] Integer number of NATS connections used
internally for geotruth queries such as object data, local vision, and floor
resolution. The default is \texttt{GOMAXPROCS}.
\item[\texttt{publishconnections}] Integer number of NATS connections used
internally to publish position updates to geotruth. The default is
\texttt{GOMAXPROCS}.
\end{description}

These values do not change the NATS connection used by a public client to call
the \texttt{pathfinding} API. They only affect internal traffic from the
pathfinding service to geotruth. Increasing them can reduce contention when many
agents query local vision and publish positions, but it also increases the
number of connections and goroutines held by the process. Values lower than
\texttt{1} are treated as one connection.

\section{Using the Service From Go}

The recommended public surface for Go integrations is formed by these packages:

\begin{itemize}
\item \texttt{embedded}
\item \texttt{pkg/pathfinding}
\item \texttt{pkg/emergency}
\item \texttt{pkg/domain}
\item \texttt{pkg/config}
\item \texttt{pkg/subjects}
\item \texttt{pkg/tuning}
\end{itemize}

\subsection{Embedded Execution}

The \texttt{embedded} package starts the same dispatcher used by the executable,
but inside the caller process. This option is useful in integrations, tests, and
tools that want explicit control over the service lifecycle.

\begin{lstlisting}[language=Go]
ctx := context.Background()
cfg := embedded.DefaultConfig()

dep, err := embedded.DefaultDependencies(ctx, cfg)
if err != nil {
    return err
}

dispatcher, err := embedded.Run(ctx, cfg, dep)
if err != nil {
    return err
}
defer dispatcher.Shutdown()
\end{lstlisting}

\texttt{DefaultConfig} returns convenient local values. In a real integration,
review at least the NATS URL, map image paths, grid dimensions, portals, and
external dependencies. \texttt{DefaultDependencies} creates internal NATS pools
for geotruth query and publishing, and uses a simulated external system. In a
real deployment, provide implementations of \texttt{domain.ExternalSystem},
\texttt{domain.GeoQuery}, and \texttt{domain.GeoPublish}. The
\texttt{GeoPublish} implementation publishes position updates and returns a
commit acknowledgement; the integrator remains responsible for releasing any
external clients or connections created outside the default pools.

To validate a local installation or integration, run:

\begin{lstlisting}
make test
\end{lstlisting}

This command runs the full suite and generates \texttt{coverage.out}, a local
per-package coverage report. To measure code exercised by all tests in the
module, including black-box and integration tests that execute other packages,
run \texttt{make test-cover-full}. The corresponding profile is
\texttt{coverage-full.out} and can be opened with \texttt{make cover-full-html}.

\subsection{Agent Pathfinding}

The \texttt{pkg/pathfinding} package provides a client over an existing NATS
connection. The client does not own the connection; the caller remains
responsible for closing it.

\begin{lstlisting}[language=Go]
pf := pathfinding.New(nc)
goal := [3]float64{10.0, 4.0, 0.0}

comm, err := pf.AgentFindPath(ctx, goal, "agent-1", 2.0, 0)
if err != nil {
    return err
}

if err := comm.MoveNCells(ctx, 5); err != nil {
    return err
}

done, err := comm.BlockingWait()
if err != nil {
    return err
}
<-done

if err := comm.ExitError(ctx); err != nil {
    return err
}
\end{lstlisting}

\texttt{AgentFindPath} starts adaptive routing toward the given goal.
\texttt{defaultCellsPerSecondMovement} defines automatic speed in cells per
second. A zero value leaves the session stopped until explicit movement commands
arrive. \texttt{moveNSteps} requests an initial number of cells to move.

To estimate the cost without moving the agent, use
\texttt{AgentNaivePathCost}:

\begin{lstlisting}[language=Go]
cost, err := pf.AgentNaivePathCost(ctx, goal, "agent-1")
\end{lstlisting}

An active session returns an \texttt{AgentCommunicator}. It can change speed,
pause, terminate, move by cells, move by meters, check whether the run is still
moving, and query the terminal error. The terminal result remains available for
a short interval after completion; after that time, or after another session for
the same agent starts, it may no longer be available.

\subsection{Subscribing to Route State}

\texttt{AgentPathSub} observes route snapshots for an agent. The channel returns
a map by floor: the key is the floor name and the value is a list of cell states
encoded with \texttt{domain.CellState}.

\begin{lstlisting}[language=Go]
states, err := pf.AgentPathSub(ctx, "agent-1")
if err != nil {
    return err
}

for snapshot := range states {
    _ = snapshot["0"]
}
\end{lstlisting}

The channel closes when the service reports that the subscription has ended.
On large maps, the service may transport one snapshot as several compressed NATS
messages; \texttt{AgentPathSub} assembles those chunks before returning the map
on the Go channel.

\subsection{Emergency Routing}

The \texttt{pkg/emergency} package controls the emergency flow field.
\texttt{Start} receives a list of exits, expressed as meter coordinates, and
optionally per-floor preference graphs.

\begin{lstlisting}[language=Go]
em := emergency.New(nc)
goals := [][3]float64{
    {1.0, 1.0, 0.0},
}

if err := em.Start(ctx, goals, nil, 2.0); err != nil {
    return err
}
defer em.Stop(ctx)
\end{lstlisting}

Flow snapshots can be observed with \texttt{FlowSub}. Each cell returns a
direction encoded with \texttt{domain.Direction}, including planar directions,
diagonal directions, and transitions between floors.

\begin{lstlisting}[language=Go]
flows, err := em.FlowSub(ctx)
if err != nil {
    return err
}

for snapshot := range flows {
    _ = snapshot["0"]
}
\end{lstlisting}

As with agent route state, \texttt{FlowSub} hides the chunked transport used when
the complete flow field is too large for one NATS message.

\subsection{Rasterized Map Tuning}

The \texttt{pkg/tuning} package helps inspect the grid generated from an image.
It does not run pathfinding, but it lets users check dimensions, scale, and wall
aura before using an image as a service map. Rasterization uses the same
production path as the service: traversable zones must be encoded in the red
channel, not only as visually light or grayscale areas.

\begin{lstlisting}[language=Go]
grid, err := tuning.TraversableImageRender(
    "./data/map.png",
    63,
    118,
    10,
    1,
)
if err != nil {
    return err
}

fmt.Println(grid)
\end{lstlisting}

For large grids, the same structure can open a temporary web visualization in
the system browser:

\begin{lstlisting}[language=Go]
if err := grid.OpenVisualizer(); err != nil {
    return err
}
\end{lstlisting}

\section{Common Workflows}

\subsection{Validate a Map Before Running It}

Before starting the service with a new image, rasterize it with the same logic
used by the service. This detects incorrect routes, wrong dimensions, walls too
close to traversable space, or maps whose traversability information is not in
the red channel.

\begin{lstlisting}[language=Go]
grid, err := tuning.TraversableImageRender(
    "./data/map.png",
    rows,
    cols,
    pxPerCell,
    wallAuraCells,
)
if err != nil {
    return err
}
fmt.Println(grid)
\end{lstlisting}

If the printed grid does not represent the space well, review
\texttt{rows}, \texttt{cols}, \texttt{pxpercell}, \texttt{cellsizem}, and
\texttt{wallauracells}. Also confirm that traversable pixels have enough red
channel value to pass the internal rasterization threshold before starting route
calculations.

\subsection{Start an Agent Route}

The basic flow requires the service to be active, NATS to be available, and the
agent to exist in geotruth. The initial position is read from geotruth and the
goal is supplied as a meter coordinate.

\begin{lstlisting}[language=Go]
pf := pathfinding.New(nc)
goal := [3]float64{x, y, z}

comm, err := pf.AgentFindPath(ctx, goal, agentID, 2.0, 0)
if err != nil {
    return err
}
\end{lstlisting}

After starting the session, wait for completion and query the final result:

\begin{lstlisting}[language=Go]
done, err := comm.BlockingWait()
if err != nil {
    return err
}
<-done

if err := comm.ExitError(ctx); err != nil {
    return err
}
\end{lstlisting}

If \texttt{ExitError} returns \texttt{nil}, the session ended without a domain
error. Otherwise, interpret the error using the error-message section.

\subsection{Control Agent Movement}

When the initial speed is zero, the session is prepared but does not advance
until it receives an explicit movement command. This mode is useful for
step-by-step simulations or for synchronizing movement with another system.

\begin{lstlisting}[language=Go]
comm, err := pf.AgentFindPath(ctx, goal, agentID, 0, 0)
if err != nil {
    return err
}

if err := comm.MoveNCells(ctx, 5); err != nil {
    return err
}
\end{lstlisting}

Movement can also be requested in meters. The call blocks until the distance is
consumed or the session ends:

\begin{lstlisting}[language=Go]
remaining, err := comm.MoveFMeters(ctx, 1.5)
if err != nil {
    return err
}
_ = remaining
\end{lstlisting}

To switch from manual mode to continuous movement, update the speed. A zero
value pauses automatic movement without destroying the session:

\begin{lstlisting}[language=Go]
_ = comm.UpdateMovementSpeed(ctx, 3.0)
_ = comm.Stop(ctx)
_ = comm.Terminate(ctx)
\end{lstlisting}

\subsection{Observe an Agent Path}

Path snapshots can be used to visualize the search state. Start the session,
subscribe to the agent identifier, and process updates until the channel closes.

\begin{lstlisting}[language=Go]
updates, err := pf.AgentPathSub(ctx, agentID)
if err != nil {
    return err
}

for snapshot := range updates {
    cellsByFloor := snapshot
    _ = cellsByFloor
}
\end{lstlisting}

Each snapshot is a map by floor. Cell values are interpreted using
\texttt{domain.CellState}.
For large maps, the Go API still yields one complete snapshot per iteration; the
split into NATS messages matters only for direct NATS consumers.

\subsection{Run an Emergency}

The emergency flow starts by defining one or more exits in meters. Optionally,
preference graphs can be supplied to favor specific routes per floor.

\begin{lstlisting}[language=Go]
em := emergency.New(nc)
goals := [][3]float64{
    {exitX, exitY, exitZ},
}

if err := em.Start(ctx, goals, routeGraphs, 2.0); err != nil {
    return err
}
\end{lstlisting}

While the emergency is active, subscribe to the flow field. The channel closes
when the service reports that publishing has ended:

\begin{lstlisting}[language=Go]
flows, err := em.FlowSub(ctx)
if err != nil {
    return err
}

for snapshot := range flows {
    directionsByFloor := snapshot
    _ = directionsByFloor
}
\end{lstlisting}

The Go channel yields complete per-floor maps. Transport chunking for large
fields is handled by the \texttt{pkg/emergency} client.

When the emergency situation ends, call \texttt{Stop} to stop computation and
remove preferred routes from the grid. The operation is idempotent: if there is
no active emergency, or if another call already started the stop, it returns
success once the system is stopped.

\begin{lstlisting}[language=Go]
if err := em.Stop(ctx); err != nil {
    return err
}
\end{lstlisting}

\subsection{Use Multiple Floors}

The previous flows also work in multi-floor spaces if the configuration includes
the corresponding layers and portals. Public coordinates are still expressed as
\texttt{[x, y, z]}; the \texttt{z} value relates the real position to geotruth
regions and grid floors.

For this case, verify that every floor has an image, portals connect traversable
cells, and the external system knows the region heights. If any of these pieces
is missing, a route that appears valid on one floor may fail when crossing to
another.

\section{NATS API}

Integrations that do not use the Go wrappers can call NATS subjects directly.
Request/reply calls return a JSON wrapper from the shared \texttt{messages}
package:

\begin{lstlisting}
{
  "ok": true,
  "error": "",
  "error_code": "",
  "data": ...
}
\end{lstlisting}

When \texttt{ok} is \texttt{false}, \texttt{error} contains the remote message
and \texttt{error\_code} can contain a stable code for known cross-process
errors. Go clients should prefer \texttt{messages.Data[T]} and
\texttt{messages.Err} to decode responses. Local wrapper errors, such as the
public operation or NATS request/response phase, are not transmitted in this
JSON; they appear only in the error returned by the Go client.

Go consumers that use NATS directly can import subject constants from
\texttt{pkg/subjects}. Consumers in other languages can use the literal subject
names shown below.

\subsection{Agent Subjects}

\begin{table}[H]
\centering
\small
\begin{tabular}{|p{0.34\textwidth}|p{0.30\textwidth}|p{0.24\textwidth}|}
\hline
\textbf{Subject} & \textbf{Input} & \textbf{Output} \\
\hline
\texttt{pathfinding.}\newline\texttt{agent\_find\_path} &
\texttt{AgentFindPath}\newline\texttt{Params} &
communicator id \\
\hline
\texttt{pathfinding.}\newline\texttt{agent\_naive\_}\newline\texttt{path\_cost} &
\texttt{AgentNaive}\newline\texttt{PathCostParams} &
cost \\
\hline
\texttt{pathfinding.}\newline\texttt{agent\_path\_watch} &
agent identifier &
stream on reply subject \\
\hline
\end{tabular}
\caption{NATS subjects for starting and observing agent pathfinding.}
\end{table}

\texttt{AgentFindPathParams} has this shape:

\begin{lstlisting}
{
  "goal": [10.0, 4.0, 0.0],
  "agentid": "agent-1",
  "defaultcellspersecondmovement": 2.0,
  "movensteps": 0
}
\end{lstlisting}

\texttt{AgentNaivePathCostParams} uses \texttt{goal} and \texttt{agentid}. The
path subscription receives the agent identifier:

\begin{lstlisting}
{"id": "..."}
\end{lstlisting}

The path subscription sends small messages to the reply subject:

\begin{lstlisting}
{
  "done": false,
  "cells": {
    "0": [0, 1, 4, 5]
  }
}
\end{lstlisting}

\subsection{Agent Communicator Subjects}

Operations on an active agent use the identifier returned by:

\begin{lstlisting}
pathfinding.agent_find_path
\end{lstlisting}

In the current implementation, this identifier is the agent identifier.

\begin{table}[H]
\centering
\small
\begin{tabular}{|p{0.34\textwidth}|p{0.30\textwidth}|p{0.24\textwidth}|}
\hline
\textbf{Subject} & \textbf{Input} & \textbf{Output} \\
\hline
\texttt{agent\_comm.}\newline\texttt{update\_movement\_speed} &
id and new speed &
OK \\
\hline
\texttt{agent\_comm.}\newline\texttt{blocking\_wait} &
id &
notification on reply subject \\
\hline
\texttt{agent\_comm.}\newline\texttt{is\_moving} &
id &
boolean \\
\hline
\texttt{agent\_comm.}\newline\texttt{exit\_error} &
id &
OK or final error \\
\hline
\texttt{agent\_comm.}\newline\texttt{terminate} &
id &
OK \\
\hline
\texttt{agent\_comm.}\newline\texttt{stop} &
id &
OK \\
\hline
\texttt{agent\_comm.}\newline\texttt{move\_n\_cells} &
id and cell count &
OK \\
\hline
\texttt{agent\_comm.}\newline\texttt{move\_f\_meters} &
id and distance in meters &
remaining meters \\
\hline
\end{tabular}
\caption{NATS subjects for controlling an agent session.}
\end{table}

The JSON bodies for these control subjects have these shapes:

\begin{lstlisting}
{"id": "..."}
{"id": "...", "float": 2.0}
{"id": "...", "n": 5}
\end{lstlisting}

The \texttt{float} field represents the new speed in
\texttt{update\_movement\_speed} and the distance in meters in
\texttt{move\_f\_meters}.

\texttt{move\_f\_meters} blocks until the requested distance is consumed or the
session ends. If the agent stopped before completing the distance, the response
can indicate \newline \texttt{ErrAgentExitedWithMetersLeft}.

\subsection{Emergency Subjects}

\begin{table}[H]
\centering
\small
\begin{tabular}{|p{0.34\textwidth}|p{0.30\textwidth}|p{0.24\textwidth}|}
\hline
\textbf{Subject} & \textbf{Input} & \textbf{Output} \\
\hline
\texttt{emergency.start} &
\texttt{EmergencyPathParams} &
OK \\
\hline
\texttt{emergency.stop} &
no body &
OK, also when no emergency was active \\
\hline
\texttt{emergency.}\newline\texttt{flow\_watch} &
no body &
stream on reply subject \\
\hline
\end{tabular}
\caption{NATS subjects for controlling emergency routing.}
\end{table}

\texttt{EmergencyPathParams} includes exits, preference graphs, and update
frequency:

\begin{lstlisting}
{
  "goals": [[1.0, 1.0, 0.0]],
  "preferencegraph": [],
  "updatetickspersecond": 2.0
}
\end{lstlisting}

\texttt{emergency.flow\_watch} messages use the same close criterion as agent
path streams. For small fields, the format is:

\begin{lstlisting}
{
  "done": false,
  "cells": {
    "0": [1, 3, 10]
  }
}
\end{lstlisting}

\subsection{Snapshot Messages in NATS Streams}

\texttt{pathfinding.agent\_path\_watch} and \texttt{emergency.flow\_watch} share
the same snapshot stream format. Go wrapper users do not need to handle this
format directly, but direct NATS consumers in other languages must accept two
valid message shapes.
The first shape is the small-message envelope preserved from the original
contract:

\begin{lstlisting}
{
  "done": false,
  "cells": {
    "0": [0, 1, 4, 5]
  }
}
\end{lstlisting}

The second shape is used when the full snapshot is too large for one NATS
message. Each message carries one compressed chunk:

\begin{lstlisting}
{
  "encoding": "gzip+json",
  "snapshotId": "42",
  "chunkIndex": 0,
  "chunkCount": 3,
  "data": "..."
}
\end{lstlisting}

Consumers must group chunks by \texttt{snapshotId}, order them by
\texttt{chunkIndex}, base64-decode each \texttt{data} field, concatenate the
resulting bytes, decompress the result with \texttt{gzip}, and decode the JSON as
a map by floor. The stream ends when this message arrives:

\begin{lstlisting}
{"done": true}
\end{lstlisting}

This chunking avoids requiring a larger NATS maximum payload. Default NATS
configurations commonly limit one message to about 1 MiB.

\section{Domain Types}

Costs use \texttt{domain.Cost}. \texttt{COST\_EMPTY\_SPACE} represents normal
movement cost and \newline \texttt{COST\_UNREACHABLE} represents an unreachable cell.

Route visualization states use \texttt{domain.CellState}:

\begin{itemize}
\item \texttt{STATE\_Unvisited}
\item \texttt{STATE\_Queued}
\item \texttt{STATE\_Expanded}
\item \texttt{STATE\_Inconsistent}
\item \texttt{STATE\_OnPath}
\item \texttt{STATE\_Obstacle}
\end{itemize}

Emergency directions use \texttt{domain.Direction}. In addition to the eight
planar directions, \texttt{DIR\_IN} and \texttt{DIR\_OUT} represent transitions
between floors, and \texttt{DIR\_UNKNOWN} means no direction is defined.

Emergency preference graphs use \texttt{domain.RouteGraph}. Each graph contains
a floor index, a list of \texttt{GraphWaypoint}, and a list of
\texttt{GraphEdge}. Waypoints use real-world meter coordinates. Graphs do not
replace the rasterized map; they modify costs to favor specific routes. If an
edge omits \texttt{Weight}, the service treats it as weight \texttt{1}.

\section{Common Error Messages}

\begin{table}[H]
\centering
\begin{tabular}{|p{0.36\textwidth}|p{0.55\textwidth}|}
\hline
\textbf{Error} & \textbf{Meaning for the user} \\
\hline
\texttt{ErrAgentCommNotFound} &
There is no active or recently completed session with the given identifier. \\
\hline
\texttt{ErrAgentNoPath} &
The service could not find a possible route to the goal. \\
\hline
\texttt{ErrAgentStillRunning} &
The final error was queried while the session was still running. \\
\hline
\texttt{ErrAgentTerminated} &
The session was explicitly terminated. \\
\hline
\texttt{ErrAgentExitedWithMetersLeft} &
The meter movement command could not be completed because the session ended
first. \\
\hline
\texttt{ErrEmergencyStillRunning} &
An emergency was started while another emergency was already active. \\
\hline
\texttt{ErrDispatcherShuttingDown} &
The service is shutting down and no longer accepts new sessions. \\
\hline
\end{tabular}
\caption{Main public API errors.}
\end{table}

Besides these domain errors, a response can fail because of NATS connectivity,
invalid JSON, coordinates outside the grid, or unavailable external geotruth
services. In those cases, review the error message together with service
configuration and dependency health.

The public Go wrappers combine errors so callers can classify the operation
without depending on message text. One error can identify the operation, the
NATS phase, and the domain cause. For example, after \texttt{Emergency.Start}, a
client can check:

\begin{lstlisting}
errors.Is(err, domain.ErrEmergencyStart)
errors.Is(err, domain.ErrNATSResponse)
errors.Is(err, domain.ErrEmergencyStillRunning)
\end{lstlisting}

The main phase sentinels are \texttt{ErrNATSRequest},
\texttt{ErrNATSResponse}, \texttt{ErrNATSSubscription}, and
\texttt{ErrNATSPublish}. Operation sentinels identify calls such as
\texttt{ErrAgentFindPath}, \texttt{ErrAgentCommStop}, or
\texttt{ErrEmergencyFlowSub}.

\subsection{Troubleshooting Guide}

When a call fails or blocks, first review these operational points:

\begin{description}
\item[NATS connection.] Check that the NATS server is active, that
\texttt{app.serverurl} matches the real address, and that the client uses the
same server. Timeouts usually indicate that the service is not subscribed to the
requested subject or that NATS is unreachable.
\item[Objects in geotruth.] Verify that the agent identifier exists in geotruth
before calling \texttt{AgentFindPath}. The service reads the agent's initial
position from that system.
\item[Images and execution directory.] Confirm that \texttt{grid.layers.img}
paths exist from the directory where the process is launched. A path that works
from the repository may fail from another directory.
\item[Grid dimensions and scale.] Review \texttt{rows}, \texttt{cols},
\texttt{pxpercell}, and \texttt{cellsizem} when the route does not match the
real space. Use \texttt{pkg/tuning} before deployment to render the grid and
detect scale or wall-rasterization problems.
\item[Cell coordinates versus rows and columns.] Internal grid coordinates are
written as \texttt{X,Y}, where \texttt{X} is the column and \texttt{Y} is the
row. Printed output and visualizers are usually read as row-column, so a cell
observed as \texttt{(row, column)} must be converted to
\texttt{(X: column, Y: row)} before multiplying by \texttt{cellsizem}.
\item[Floors and regions.] If a coordinate \texttt{[x, y, z]} does not resolve
to the expected floor, review the external system that translates real points
to regions and the names defined in \texttt{grid.layers}. Floor names must be
consistent between service configuration and geotruth.
\item[Portals between floors.] In multi-floor environments, check that each
portal joins traversable cells, source and destination layer indexes exist, and
the traversal cost is positive. A badly placed portal can make an inter-floor
route impossible even if each floor is traversable by itself.
\item[Route not found.] For \texttt{ErrAgentNoPath}, inspect whether the exit or
goal is inside a blocked zone, dynamic obstacles closed the only available
connection, or the rasterized grid isolated the initial cell.
\item[Blocking calls.] \texttt{BlockingWait} and \texttt{MoveFMeters} can wait
until the session ends or the agent consumes the requested distance. If many
simultaneous calls are made for different agents, review
\texttt{pathfinding.agent.blockingwaithandlerworkers} and
\texttt{pathfinding.agent.movefmetershandlerworkers}. These values only limit
NATS handlers; they do not change route calculation or command ordering for a
single agent.
\end{description}

\end{document}
